\section{The intrinsic inverse and the completion defect}
\label{sec:inverse}
This section records the deterministic inverse used after the local gates
have been constructed. It preserves the information category of the
preceding manuscript: a finite contact fingerprint recognizes a branch;
two density \emph{functions}, not that fingerprint, reconstruct its geometry.
The full earlier proofs and their broader variants remain in Supplement Q.

\subsection{Action and curvature from two densities}
Let $x(s),y(u)$ parametrize the facing arcs by arclength and put
$D(s,u)=|x(s)-y(u)|$. The stationary two-flight action is
\[
 W(s,t)=D(s,u(s,t))+D(t,u(s,t)),\qquad
                       D_u(s,u)+D_u(t,u)=0.
\]
The intermediate second derivative is positive near the normal orbit.
In section coordinates the incoming conjugate momentum is $-W_s$;
Liouville measure consequently has factor $(-W_{st})\dd s\dd t\dd r$,
where $r$ is residual time before the first impact. In a collar shorter
than adjacent free flights the admissible residual interval is
$0<r<T-W$. Therefore
\begin{equation}\label{eq:law}
 f_T(s,t)=\frac{-W_{st}(s,t)}{2\pi A}(T-W(s,t))_+.
\end{equation}
Channel clearance and the strict action minimum make this precisely the
selected three-impact event. There is no conditioning on success.
On a strictly active rectangle, subtraction at two times gives
\begin{equation}\label{eq:ratio}
 W=\frac{T_1f_2-T_2f_1}{f_2-f_1},\qquad
 A=\frac{(T_2-T_1)(-W_{st})}{2\pi(f_2-f_1)}.
\end{equation}
In particular the absolute normalization, and not merely the conditional
endpoint distribution, determines $A$.

On the diagonal write
\[
 \ell=\tfrac12W(s,s),\quad a=W_s(s,s),\quad v=(1-a^2)^{1/2}.
\]
If $T,N$ are the source tangent and inward normal, the vector from the
opposite nearest point to the source is $e=aT+vN$. At that point
\[
 D_{ss}=v^2/\ell+\kappa v,\quad D_{su}=-v/\ell,
 \quad D_{uu}=1/\ell+\kappa_o.
\]
The stationary Schur complement gives
\[
 W_{ss}=D_{ss}-\frac{D_{su}^2}{2D_{uu}},\qquad
 W_{st}=-\frac{D_{su}^2}{2D_{uu}}.
\]
Hence
\begin{align}
 \kappa(s)&=\frac{W_{ss}-W_{st}-v^2/\ell}{v},\label{eq:curv-source}\\
 \kappa_o(u(s))&=\ell^{-1}
                   \left\{\frac{v^2}{-2\ell W_{st}}-1\right\},
 &u'(s)&=\frac{v}{1+\ell\kappa_o(u(s))}>0.\label{eq:curv-target}
\end{align}
All action derivatives here are evaluated at $(s,s)$. Integrating the
source Frenet equations reconstructs $x$. The opposite arc is
$y(u(s))=x(s)-\ell(aT+vN)$. The gap is $W(0,0)/2$.
This reconstructs smooth functions on intervals, not only Taylor series.

\begin{proposition}[Conditional local inverse]
\label{prop:inverse}
Assume physical $C^3$ bounds on larger contact arcs; positive gap, area,
curvature, clearance, adjacent-flight, twist, active and time-separation
margins; $v\ge v_*>0$; and uniformly bounded $C^2$ densities with
$f_2-f_1$ bounded below. For two physical density pairs differing by
at most $e$ in $C^2$, the gaps, free areas and both intrinsic $C^2$
contact arcs differ by at most $Ce$, modulo their common reflection.
\end{proposition}
\begin{proof}
Division in \eqref{eq:ratio} is locally Lipschitz in $C^2$ with the
stated denominator margin. The curvature formulas thus give uniform
curvature errors $Ce$. Continuous dependence in the Frenet system gives
the source-arc estimate. For the opposite arc integrate the positive
speed in \eqref{eq:curv-target}, compare its inverse parameters, and
use the physical $C^3$ bound to make the opposite curvature uniformly
Lipschitz in its own arclength. A second Frenet integration gives its
$C^2$ estimate. This avoids differentiating the foot-point formula twice
and silently demanding a third derivative of $W$. The area follows
from \eqref{eq:ratio}.
\end{proof}

The local canonical relation has momenta $p=-W_s$, $p'=W_t$.
Knowing it determines $W$ only up to an additive constant on the local
branch; an absolute action value and $A$ are additionally needed to
recover the densities. Conversely \eqref{eq:ratio} determines all of
these quantities. This is the precise relation to local travel-time and
lens data, not an identification with a marked length spectrum.

\subsection{Whole analytic images and assembly}
For fitted physical density pairs, Proposition~\ref{prop:inverse} gives
$Ce$ errors on arcs. Positive curvature converts them to support-function
errors on a fixed normal-angle interval, using the $C^3$ bounds for
parameter inversion. Slit a narrower periodic analytic strip along a
shorter such interval. The harmonic measure of the slit is positive and
has a positive minimum $\theta$ on a further narrower closed strip.
The subharmonic maximum principle for the logarithm of the support
 difference gives an $O(e^\theta)$ bound there; Cauchy estimates give
the same bound through two derivatives on the real circle. Slit endpoints
are regular, and zeros are treated by regularization. Thus the entire
pair images have error
\begin{equation}\label{eq:Delta}
                         \Delta=Ce^\theta,\qquad 0<\theta<1.
\end{equation}
The exponent depends on the analytic prior and continuation geometry,
as in the conditional-continuation setting of \cite{Trefethen}.

Cluster centered complete images modulo $O(2)$. Same-type estimates have
distance $O(\Delta)$, whereas different types have distance at least
$\sigma-O(\Delta)$. Threshold $\sigma/2$ therefore gives exactly the
visible types for small enough $\Delta$. Each record is one multigraph
edge; loops and repetitions are retained. Choose a spanning tree in a
connected component and place its representatives using least-discrepancy
matches. The symmetry margins \eqref{eq:rot}--\eqref{eq:refl} force any
two accurate matches to have the same reflection component and rotation
angles differing by $O(\Delta/\eta)$. Steiner points control translations.
Tree paths have at most $r_0-1$ edges, so all placement errors are
$O(\Delta)$, independent of the number of repeated observations.
For each non-tree edge the translated target gives a cycle displacement
$d_e\in\Lambda$, with error $O(\Delta)$ and norm at most $B_0$.
No supplied integer edge label is used.

Let $I$ denote the visible types in the component and set
$\Gamma=\sum_e\Z d_e$. If its rank is two, put
\begin{equation}\label{eq:defect}
 \mathcal D=\covol\Gamma-A-\sum_{i\in I}\area(C_i)
 =([\Lambda:\Gamma]-1)\covol\Lambda
                         +\sum_{i\notin I}\area(C_i).
\end{equation}
Both terms on the right are nonnegative. Thus $\mathcal D=0$ exactly
when all types are visible and $\Gamma=\Lambda$. Otherwise
$\mathcal D\ge g_0:=\min(V_0,a_0)$. A rank-deficient component does
not certify completion. Omitting or repeating genuine records cannot
make an incomplete component pass the exact test.

\subsection{Recovering integer structure before the area test}
Suppose all cycle-vector errors are at most $t$. The determinant error
of every pair is at most $(2B_0+t)t$. A nonzero true determinant is
an integer multiple of $\covol\Lambda$, hence at least $V_0$.
For sufficiently small $t$, threshold $V_0/2$ distinguishes rank two,
and maximizing the estimated determinant selects an independent true
pair $D=[d_1\ d_2]$. It satisfies
\[
 V_0\le|\det D|\le B_0^2,\qquad
 n_D=|\det D|/\covol\Lambda\in\{1,\ldots,Q\}.
\]
Every coordinate of $D^{-1}d_e$ lies in $n_D^{-1}\Z$ and has bounded
absolute value. Distinct reduced rationals of denominators at most $Q$
are separated by at least $Q^{-2}$. Uniform matrix-inversion bounds
therefore let a finite rational search recover all these coordinates
exactly once their errors are less than $1/(4Q^2)$.

Let $q$ be their least common denominator. All denominators divide
$n_D$, so $q\le Q$. The integer column group generated by $qI_2$
and the vectors $qD^{-1}d_e$ contains $q\Z^2$. A column Hermite basis
$H$ gives
\begin{equation}\label{eq:Hermite}
 \Gamma=(D/q)H\Z^2,\qquad
 \covol\Gamma=|\det D|\,|\det H|/q^2.
\end{equation}
There are finitely many possible $H$ for $q\le Q$; in particular
$1\le|\det H|\le q^2$. Once the rational coordinates are fixed, the
covolume error is $O(t)$. This is not integer spanning of noisy real
vectors, which need not even generate a discrete group.

For a bounded support function,
$\area(C)=\tfrac12\int_0^{2\pi}(p^2-(p')^2)\dd\varphi$ is locally
Lipschitz in $C^1$. Use one area estimate per visible type, one fitted
estimate of the common free area, and \eqref{eq:Hermite} to estimate
\eqref{eq:defect}. Its error is $O(\Delta)$.

\begin{theorem}[Uniform finite-list certificate]\label{thm:certificate}
There are $C,e_*>0$ and $0<\theta<1$, depending only on the uniform
class, such that a finite list of genuine records with simultaneous
$C^2$ density error $e<e_*$ permits correct visible incidence and
rational period recovery, and defect estimates of error at most
$Ce^\theta<g_0/4$. Accept a rank-two component precisely when
$|\widehat{\mathcal D}|<g_0/2$. Every accepted component is complete
and saturated, and its whole-table error is at most $Ce^\theta$.
Every complete saturated component passes. The same event controls all
components, repetitions and data-dependent deletions without a graph
union bound. No witness or reference table is input to the test.
\end{theorem}
\begin{proof}
The preceding estimates recover the exact clusters and finite rational
structure. The gap following \eqref{eq:defect} separates the two cases.
On acceptance the estimated period basis is $\widehat D H/q$, compared
with the true basis $DH/q$. The finite set of allowed Hermite forms
bounds their error by $Ce^\theta$, simultaneously with the placed-body
errors. One may project to the compact physical class at an additional
constant factor. Physical fitting and projection can be defined using
near-minimizers in fixed countable dense families and a first-eligible
selection. Their observables are continuous on the protected physical
class, so this yields measurable choices. This is an existence inverse,
not an efficient search over all analytic bodies.
\end{proof}
